\documentclass[11pt,a4paper]{ctexart}

\usepackage[margin=2.0cm]{geometry}
\usepackage{amsmath}
\usepackage{xcolor}
\usepackage{tikz}
\usepackage[hidelinks]{hyperref}

\usetikzlibrary{arrows.meta,positioning,calc}

\setlength{\parindent}{2em}
\setlength{\parskip}{0.45em}
\linespread{1.15}
\sloppy

\definecolor{lotblue}{HTML}{2563EB}
\definecolor{lotgreen}{HTML}{16A34A}
\definecolor{lotorange}{HTML}{EA580C}
\definecolor{lotred}{HTML}{DC2626}
\definecolor{lotgray}{HTML}{64748B}
\definecolor{lotline}{HTML}{334155}
\definecolor{lotfill}{HTML}{F8FAFC}
\definecolor{lotsoftblue}{HTML}{DBEAFE}
\definecolor{lotsoftgreen}{HTML}{DCFCE7}
\definecolor{lotsoftorange}{HTML}{FFEDD5}
\definecolor{lotsoftred}{HTML}{FEE2E2}
\definecolor{lotsoftgray}{HTML}{E2E8F0}

\newcommand{\code}[1]{\texttt{\detokenize{#1}}}

\tikzset{
  lot/.style={
    draw=lotline,
    rounded corners=2pt,
    fill=lotfill,
    align=center,
    minimum width=2.0cm,
    minimum height=1.05cm,
    inner xsep=4pt,
    inner ysep=4pt,
    font=\small
  },
  emptylot/.style={lot, draw=lotgray, fill=lotsoftgray},
  openlot/.style={lot, draw=lotblue, fill=lotsoftblue},
  closelot/.style={lot, draw=lotgreen, fill=lotsoftgreen},
  fulllot/.style={lot, draw=lotorange, fill=lotsoftorange},
  stoplot/.style={lot, draw=lotred, fill=lotsoftred},
  arrow/.style={-{Latex[length=2.3mm]}, thick, draw=lotline},
  bluearrow/.style={-{Latex[length=2.3mm]}, thick, draw=lotblue},
  greenarrow/.style={-{Latex[length=2.3mm]}, thick, draw=lotgreen},
  orangearrow/.style={-{Latex[length=2.3mm]}, thick, draw=lotorange},
  redarrow/.style={-{Latex[length=2.3mm]}, thick, draw=lotred},
  note/.style={align=left, font=\small, text=lotline}
}

\title{Multi-Lot 操作机制说明}
\author{glftmm\_lot}
\date{2026-06-16}

\begin{document}
\maketitle

\section{Lot 栈定义}

当前只看 lot 的操作机制，不重复策略 overview。

现在一共有 5 个 lot。起始状态下，5 个 lot 全部都是空 lot。
从内层到外层，依次标记为 lot 0、lot 1、lot 2、lot 3、lot 4。

\begin{center}
\begin{tikzpicture}[node distance=5mm]
  \node[emptylot] (l0) {lot 0\\空};
  \node[emptylot, right=of l0] (l1) {lot 1\\空};
  \node[emptylot, right=of l1] (l2) {lot 2\\空};
  \node[emptylot, right=of l2] (l3) {lot 3\\空};
  \node[emptylot, right=of l3] (l4) {lot 4\\空};
  \node[note, above=7mm of l0, align=center] {内层};
  \node[note, above=7mm of l4, align=center] {外层};
  \node[note, below=6mm of l2, text width=11.5cm] {
    初始时没有任何 lot 持仓。第一层仓位只能先从 lot 0 开始建立；
    后面的 lot 不会提前启用。
  };
\end{tikzpicture}
\end{center}

\section{Open / Close 权限}

核心规则只有一句：只有最外层的非空 lot 能正常 open + close。

如果某一层的外层已经是非空 lot，这一层就只能 close-only。
close-only 的意思是只能减仓或平仓，不能继续 open。
它要么等到水上平仓离场，要么触发 stoploss 止损离场。

\begin{center}
\begin{tikzpicture}[node distance=5mm]
  \node[closelot] (l0) {lot 0\\有仓位\\close-only};
  \node[closelot, right=of l0] (l1) {lot 1\\有仓位\\close-only};
  \node[openlot, right=of l1] (l2) {lot 2\\有仓位\\open + close};
  \node[emptylot, right=of l2] (l3) {lot 3\\空\\等待};
  \node[emptylot, right=of l3] (l4) {lot 4\\空\\等待};

  \draw[bluearrow] ($(l2.north)+(0,0.55)$) -- ($(l2.north)+(0,0.05)$);
  \node[note, above=8mm of l2, align=center] {最外层非空 lot\\正常 open + close};

  \node[note, below=6mm of l2, text width=12cm] {
    这里 lot 2 是最外层非空 lot，所以 lot 2 能 open + close。
    lot 0 和 lot 1 的外层已经有非空 lot，所以它们只能 close-only。
  };
\end{tikzpicture}
\end{center}

\section{下一层开启条件}

只有第 \(n\) 层满仓了，才会开启第 \(n+1\) 层。
一旦第 \(n+1\) 层开启，第 \(n\) 层进入 close-only，不再接新的 open。

\begin{center}
\begin{tikzpicture}
  \node[openlot] (a0) at (0,0) {lot n\\非空未满\\open + close};
  \node[emptylot] (a1) at (5.0,0) {lot n+1\\空\\等待};

  \node[fulllot] (b0) at (0,-2.4) {lot n\\已满仓\\转 close-only};
  \node[openlot] (b1) at (5.0,-2.4) {lot n+1\\开启\\open + close};

  \draw[orangearrow] (a0.south) -- node[left,font=\small] {lot n 满仓} (b0.north);
  \draw[bluearrow] (a1.south) -- node[right,font=\small] {被开启} (b1.north);
  \draw[bluearrow] (b0.east) -- node[below,font=\small] {新 open 交给外层} (b1.west);
\end{tikzpicture}
\end{center}

例如 lot 0 满仓后，lot 1 才开启；此时 lot 0 进入 close-only。
如果 lot 1 也满仓，lot 2 才开启；此时 lot 1 进入 close-only。

\section{订单处理顺序}

一个订单来了，先从内层开始看能不能关掉内层仓位。
顺序固定从里到外刷一遍：

\[
\text{lot 0} \rightarrow \text{lot 1} \rightarrow \text{lot 2}
\rightarrow \text{lot 3} \rightarrow \text{lot 4}
\]

如果 lot 0 可以被这个订单 close，就先 close lot 0。
如果订单还有剩余量，再看 lot 1；再剩余，再看 lot 2。
这个过程一直刷到最外层。

\begin{center}
  \begin{tikzpicture}[node distance=10mm]
  \node[closelot] (l0) {lot 0\\先看};
  \node[closelot, right=of l0] (l1) {lot 1\\再看};
  \node[openlot, right=of l1] (l2) {lot 2\\再看};
  \node[emptylot, right=of l2] (l3) {lot 3\\再看};
  \node[emptylot, right=of l3] (l4) {lot 4\\最后};

  \draw[greenarrow] (l0) -- node[above,font=\scriptsize] {剩余} (l1);
  \draw[greenarrow] (l1) -- node[above,font=\scriptsize] {剩余} (l2);
  \draw[greenarrow] (l2) -- node[above,font=\scriptsize] {剩余} (l3);
  \draw[greenarrow] (l3) -- node[above,font=\scriptsize] {剩余} (l4);

  \node[note, below=6mm of l2, text width=12cm] {
    这一步的重点是：订单优先给内层退出机会。
    即使最外层非空 lot 可以 open + close，撮合检查仍然从 lot 0 开始。
  };
\end{tikzpicture}
\end{center}

\section{内部层平完后的旋转}

如果内部某一层被平完了，这一层会被挪到最外层。
它外面的层整体向内移动一格，所以这些外部层的 index 都减 1。

例子：lot 0、lot 1、lot 2 都有仓位，其中 lot 2 是最外层非空 lot。
如果内部的 lot 1 被平完，lot 1 会挪到最外层；
原来的 lot 2 向内移动，index 从 2 变成 1。

\begin{center}
\begin{tikzpicture}[node distance=5mm]
  \node[note] at (-0.2,1.25) {平完前};
  \node[closelot] (a0) at (0,0) {lot 0\\有仓位};
  \node[closelot] (a1) at (2.25,0) {lot 1\\平完};
  \node[openlot] (a2) at (4.50,0) {lot 2\\有仓位};
  \node[emptylot] (a3) at (6.75,0) {lot 3\\空};
  \node[emptylot] (a4) at (9.00,0) {lot 4\\空};

  \node[note] at (-0.2,-1.65) {平完后};
  \node[closelot] (b0) at (0,-2.9) {lot 0\\有仓位};
  \node[openlot] (b1) at (2.25,-2.9) {lot 1\\原 lot 2};
  \node[emptylot] (b2) at (4.50,-2.9) {lot 2\\原 lot 3};
  \node[emptylot] (b3) at (6.75,-2.9) {lot 3\\原 lot 4};
  \node[emptylot] (b4) at (9.00,-2.9) {lot 4\\原 lot 1};

  \draw[arrow] (a0) -- (b0);
  \draw[bluearrow] (a2) -- node[right,font=\scriptsize] {index - 1} (b1);
  \draw[bluearrow] (a3) -- node[right,font=\scriptsize] {index - 1} (b2);
  \draw[bluearrow] (a4) -- node[right,font=\scriptsize] {index - 1} (b3);
  \draw[orangearrow] (a1) .. controls (5.0,-1.0) and (8.5,-1.6) .. node[right,font=\scriptsize] {挪到最外层} (b4);
\end{tikzpicture}
\end{center}

\section{七条规则}

\begin{enumerate}
  \item 当前是 5 个 lot。
  \item 起始状态是 5 个 lot 均为空 lot。
  \item 5 个 lot 从内到外依次标记为 0、1、2、3、4。
  \item 只有最外层的非空 lot 能正常 open + close；如果一层的外层是非空的，这一层只能 close-only，要么等水上平仓离场，要么止损离场。
  \item 只有第 \(n\) 层满仓了，才能开启第 \(n+1\) 层；此时第 \(n\) 层进入 close-only。
  \item 一个订单来了，优先从内层开始看是否能够关掉内层仓位，从里到外刷一遍。
  \item 当内部某一层被平完了，把这一层挪到最外层，其外部层的 index 减 1。
\end{enumerate}

\end{document}
